\documentclass[reqno,11pt]{amsart}
\usepackage{math327-lecture}
\usepackage{needspace}
\lectureinfo{11}{Sep. 28, 2026}{Equivalence Relations and Lagrange's Theorem}

% Continues Lecture 10; material moved from the commented-out end of lecture-10.tex.
\begin{document}
\makelecturetitle


\section{Recollections: cosets form a partition}

Last time we introduced the notion of left and right cosets of a group. These were defined as follows: given a group $G$ and a subgroup $H \leq G$, a left coset is a set of the form \be g H = \{ g h \mid h \in H \} \qquad \text{ for some fixed } g \in G \ee and a right coset is one of the form \be{} H g = \{ hg \mid h \in H \} \qquad \text{ for some fixed } g \in G. \ee

We ended last class by proving the following result:

\begin{lemma}[Equality of cosets]\label{lem:coset-equality}
	Let $H\leq G$ and $g_1,g_2\in G$. The following are equivalent:
	\begin{enumerate}[label=\textbf{(\arabic*)}]
		\item $g_1H\cap g_2H\neq\varnothing$;
		\item $g_1H=g_2H$;
		\item $g_2^{-1}g_1\in H$;
		\item $g_1^{-1}g_2\in H$.
	\end{enumerate}
\end{lemma}
%\begin{proof}
%	First consider the special case $g_2=e_G$. We show that
%	\[
%	gH\cap H\neq\varnothing
%	\quad\Longleftrightarrow\quad g\in H
%	\quad\Longleftrightarrow\quad gH=H.
%	\]
%	We prove $g\in H\Rightarrow gH=H\Rightarrow gH\cap H\neq\varnothing
%	\Rightarrow g\in H$.
%	\begin{itemize}
%		\item If $g\in H$, closure gives $gH\subseteq H$. Conversely, every
%		$h\in H$ can be written $h=g(g^{-1}h)$, with $g^{-1}h\in H$, so
%		$H\subseteq gH$. Thus $gH=H$.
%		\item If $gH=H$, then $gH\cap H=H$, which is nonempty since
%		$e_G\in H$.
%		\item If $z\in gH\cap H$, write $z=gh$ with $h\in H$. Then
%		$g=zh^{-1}\in H$.
%	\end{itemize}
%	This completes the special case.
%	
%	For general $g_1,g_2$, left multiplication by $g_2^{-1}$ is a bijection
%	of $G$. Hence
%	\begin{align*}
%		g_1H\cap g_2H\neq\varnothing
%		&\quad\Longleftrightarrow\quad
%		(g_2^{-1}g_1)H\cap H\neq\varnothing\\
%		&\quad\Longleftrightarrow\quad g_2^{-1}g_1\in H\\
%		&\quad\Longleftrightarrow\quad (g_2^{-1}g_1)H=H\\
%		&\quad\Longleftrightarrow\quad g_1H=g_2H.
%	\end{align*}
%	Finally, $g_2^{-1}g_1\in H$ if and only if its inverse
%	$g_1^{-1}g_2$ belongs to $H$.
%\end{proof}

Recall the following definition: 

\begin{define}[Partition]
	A \emph{partition} of a set $X$ is a collection of nonempty subsets
	$\{X_\alpha\}_{\alpha\in I}$ such that
	\[
	X=\bigcup_{\alpha\in I}X_\alpha,
	\qquad X_\alpha\cap X_\beta=\varnothing\quad(\alpha\neq\beta).
	\]
	We write $X=\bigsqcup_{\alpha\in I}X_\alpha$.
\end{define}

Now, we can prove the basic fact: 

\begin{proposition}
	The distinct left cosets of $H$ form a partition of $G$.
\end{proposition}
\begin{proof}
	Every $g\in G$ belongs to $gH$, since $g=ge_G$ and $e_G\in H$.
	Thus the cosets are nonempty and cover $G$. By
	Lemma~\ref{lem:coset-equality}, two cosets that intersect are equal, so
	distinct cosets are disjoint.
\end{proof}



\section{Example: cosets of \texorpdfstring{$8\Z$}{8Z}}\label{sec:8Z}

In the additive group $G=\Z$, take $H=8\Z$. Cosets are written additively:
\begin{equation}
	a+8\Z=\{a+8k:k\in\Z\}.
\end{equation}
For instance,
\begin{align*}
	0+8\Z&=\{\ldots,-16,-8,0,8,16,\ldots\},\\
	1+8\Z&=\{\ldots,-15,-7,1,9,17,\ldots\},\\
	7+8\Z&=\{\ldots,-9,-1,7,15,23,\ldots\}.
\end{align*}
Arranging the integers by their remainders gives eight disjoint columns:
\[
\begin{array}{c|rrrrrrrr}
	\text{remainder}&0&1&2&3&4&5&6&7\\ \hline
	&\vdots&\vdots&\vdots&\vdots&\vdots&\vdots&\vdots&\vdots\\
	&16&17&18&19&20&21&22&23\\
	&8&9&10&11&12&13&14&15\\
	&0&1&2&3&4&5&6&7\\
	&-8&-7&-6&-5&-4&-3&-2&-1\\
	&-16&-15&-14&-13&-12&-11&-10&-9\\
	&\vdots&\vdots&\vdots&\vdots&\vdots&\vdots&\vdots&\vdots
\end{array}
\]
Each column is one coset. The equality criterion becomes
\begin{equation}
	a+8\Z=b+8\Z
	\quad\Longleftrightarrow\quad a-b\in8\Z
	\quad\Longleftrightarrow\quad a\equiv b\pmod8.
\end{equation}
For example, $0+8\Z=8+8\Z$ and $1+8\Z=17+8\Z$.
Every integer has exactly one remainder in $\{0,1,\ldots,7\}$, so
\[
\Z=\bigsqcup_{r=0}^7(r+8\Z),\qquad
\Z/8\Z=\{0+8\Z,1+8\Z,\ldots,7+8\Z\}.
\]
Writing $\overline a=a+8\Z$, we may abbreviate this as
\[
\Z/8\Z=\{\overline0,\overline1,\ldots,\overline7\}.
\]
Thus an infinite group can have a finite set of cosets.

%\section{Optional: equivalence relations and partitions}
%
%An equivalence relation is a way to declare elements of a set equivalent
%while retaining the basic properties of equality.
%
%\begin{define}[Equivalence relation]
%Let $X$ be a set. A relation $\sim$ on $X$ is an \emph{equivalence relation}
%if, for all $a,b,c\in X$, it satisfies:
%\begin{enumerate}[label=\textbf{(\arabic*)}]
%  \item \emph{Reflexivity:} $a\sim a$.
%  \item \emph{Symmetry:} if $a\sim b$, then $b\sim a$.
%  \item \emph{Transitivity:} if $a\sim b$ and $b\sim c$, then $a\sim c$.
%\end{enumerate}
%The \emph{equivalence class} of $a\in X$ is
%\[
%  C_a=\{z\in X:z\sim a\}.
%\]
%\end{define}
%
%\begin{proposition}
%The distinct equivalence classes form a partition of $X$. More precisely,
%\[
%  C_a\cap C_b\neq\varnothing
%  \quad\Longleftrightarrow\quad C_a=C_b
%  \quad\Longleftrightarrow\quad a\sim b.
%\]
%\end{proposition}
%\begin{proof}
%Reflexivity gives $a\in C_a$, so the classes are nonempty and cover $X$.
%If $z\in C_a\cap C_b$, then $z\sim a$ and $z\sim b$. Symmetry and
%transitivity give $a\sim b$. If $a\sim b$ and $w\in C_a$, then
%$w\sim a\sim b$, so $w\in C_b$. This gives $C_a\subseteq C_b$; symmetry
%gives the reverse inclusion. Finally, if $C_a=C_b$, their intersection
%contains $a$. Thus classes are either equal or disjoint.
%\end{proof}
%
%Choose a set $R\subseteq X$ containing exactly one representative of each
%class. Then
%\[
%  X=\bigsqcup_{a\in R}C_a.
%\]
%Conversely, any partition of $X$ defines an equivalence relation: declare
%$a\sim b$ when $a$ and $b$ lie in the same part.
%
%\begin{example}[Congruence modulo eight]
%On $\Z$, define $a\sim b$ if $8\mid(b-a)$. This is an equivalence
%relation (transitivity uses $c-a=(c-b)+(b-a)$), and its class
%$C_a=a+8\Z$ is exactly the column containing $a$ in the table of cosets of $8\Z$
%from Lecture~10.
%Besides $R=\{0,1,\ldots,7\}$ we could use $R'=\{-3,-2,\ldots,4\}$:
%\[
%  \Z=\bigsqcup_{a\in R}C_a=\bigsqcup_{a\in R'}C_a.
%\]
%The classes are fixed by the relation; the representatives are not.
%\end{example}
%
%
%Let $H\leq G$. Define a relation on $G$ by
%\begin{equation}
%  a\sim b\quad\Longleftrightarrow\quad b^{-1}a\in H.
%\end{equation}
%It is an equivalence relation:
%\begin{enumerate}[label=\textbf{(\arabic*)}]
%  \item $a^{-1}a=e_G\in H$, so $a\sim a$.
%  \item If $b^{-1}a\in H$, then $(b^{-1}a)^{-1}=a^{-1}b\in H$, so
%  $b\sim a$.
%  \item If $b^{-1}a\in H$ and $c^{-1}b\in H$, then
%  \[
%    c^{-1}a=(c^{-1}b)(b^{-1}a)\in H,
%  \]
%  so $a\sim c$.
%\end{enumerate}
%The class of $a$ is precisely its left coset:
%\begin{equation}
%  C_a=\{g\in G:a^{-1}g\in H\}=aH.
%\end{equation}
%Indeed, $a^{-1}g=h\in H$ if and only if $g=ah$.
%This gives another explanation of why left cosets partition $G$.
%
%\begin{remark}[The relation for right cosets]
%To obtain right cosets, use the relation
%\[
%  a\sim_R b\quad\Longleftrightarrow\quad ab^{-1}\in H.
%\]
%Then the class of $a$ is $\{g\in G:ga^{-1}\in H\}=Ha$.
%By contrast, replacing $b^{-1}a\in H$ with $a^{-1}b\in H$ does
%\emph{not} change the left-coset relation: these two elements are inverses,
%and $H$ is closed under inverses.
%\end{remark}


\section{$G/H$ and Lagrange's theorem}

For a subgroup $H \leq G$, recall that 

\begin{define} We write $G/H$ for the set of left cosets of $G$. So every element of $G/H$ is a coset of the form $gH$, and we often denote this as $\overline{g}$. Note that we are only recording the \textit{distinct} left cosets here, and it could very well be that $g_1H  =g_2H$, i.e. $\overline{g}_1 = \overline{g}_2$. 
	If $R\subseteq G$ contains exactly one representative of each left coset,
	the precise formulas are
	\begin{equation}
		G/H=\{rH:r\in R\},\qquad G=\bigsqcup_{r\in R}rH.
	\end{equation}
	
	
	From our Lemma above, this happens whenever $g_1^{-1} g_2\in H$ or $g_2^{-1}g_1 \in H.$ \end{define} 

One of the first questions we can ask is: how big is $G/H$? The answer is recorded in a quantity called the index of $H$ in $G$.

\begin{define}[Index]
	The \emph{index} of a subgroup $H$ in $G$, denoted $[G:H]$, is the number
	of distinct left cosets:
	\[
	[G:H]=|G/H|.
	\]
	The index may be infinite.
\end{define}


\begin{example}[Indices in $S_3$ and $\Z$]
\leavevmode
\begin{enumerate}[label=\textbf{(\alph*)}]
  \item \emph{Subgroups of $S_3$.} As in Lecture~10, write
  $S_3=\{1,x,x^2,y,xy,x^2y\}$ with $x^3=1$, $y^2=1$, $yx=x^2y$.
  \begin{itemize}
    \item For $N=\{1,x,x^2\}$ we found $S_3/N=\{N,yN\}$, so $[S_3:N]=2$.
    \item For $H=\{1,y\}$ we found $S_3/H=\{H,xH,x^2H\}$, so
    $[S_3:H]=3$.
    \item For the trivial subgroup $I=\{1\}$, each coset $gI=\{g\}$ is a
    single element, and distinct elements give distinct cosets. Thus
    $S_3/I=\{\{g\}:g\in S_3\}$ and $[S_3:I]=6$.
    \item At the other extreme, $gS_3=S_3$ for every $g$, so
    $[S_3:S_3]=1$.
  \end{itemize}
  In each case $|S_3|=6$ equals $[S_3:K]\cdot|K|$:
  \[
    6=2\cdot3=3\cdot2=6\cdot1=1\cdot6.
  \]

  \item \emph{The subgroups $n\Z\leq\Z$.} Let $n\geq1$. By the division
  algorithm, every $a\in\Z$ can be written $a=qn+r$ with
  $0\leq r\leq n-1$, so $a-r\in n\Z$ and $a+n\Z=r+n\Z$. Hence every
  coset is one of
  \[
    0+n\Z,\ 1+n\Z,\ \ldots,\ (n-1)+n\Z.
  \]
  These are distinct: if $0\leq r<s\leq n-1$, then $0<s-r<n$, so
  $s-r\notin n\Z$ and $r+n\Z\neq s+n\Z$. Therefore
  \[
    \Z/n\Z=\{\overline0,\overline1,\ldots,\overline{n-1}\},
    \qquad [\Z:n\Z]=n.
  \]
  The case $n=8$ is the table in Section~\ref{sec:8Z}. Although $\Z$ and
  $n\Z$ are both infinite, the index is finite.

  For $n=0$, the subgroup is $0\Z=\{0\}$, each coset $a+0\Z=\{a\}$ is a
  single integer, and $[\Z:0\Z]$ is infinite.
\end{enumerate}
\end{example}






\begin{theorem}[Lagrange]
If $G$ is finite and $H\leq G$, then
\begin{equation}
  |G|=[G:H]\,|H|.
\end{equation}
In particular, $|H|$ divides $|G|$.
\end{theorem}
\begin{proof}
Choose one representative of each left coset, say $g_1,\ldots,g_m$, where
$m=[G:H]$. Since the cosets partition $G$ and each has $|H|$ elements,
\[
  G=\bigsqcup_{i=1}^m g_iH,
  \qquad |G|=\sum_{i=1}^m|g_iH|=m|H|.
\]
\end{proof}


\begin{lemma}
	For every subgroup $H\leq G$ and every $g\in G$, the map
	\[
	H\longrightarrow gH,\qquad h\longmapsto gh,
	\]
	is a bijection.
\end{lemma}
\begin{proof}
	It is surjective by the definition of $gH$. If $gh_1=gh_2$, cancellation
	gives $h_1=h_2$, so it is injective. Its inverse sends $z\in gH$ to
	$g^{-1}z\in H$.
\end{proof}

More generally, left multiplication by $g$ is a bijection $G\to G$.
In particular, if $H$ is finite, every coset has the same number of
elements as $H$.


\begin{corollary}
If $G$ is finite and $a\in G$, then the order of $a$ divides $|G|$.
\end{corollary}
\begin{proof}
Apply Lagrange's theorem to the cyclic subgroup $\langle a\rangle$, whose
cardinality is the order of $a$.
\end{proof}

\begin{corollary}
Every group of prime order $p$ is cyclic, hence isomorphic to the cyclic
group $C_p$ of order $p$.
\end{corollary}
\begin{proof}
Choose $a\neq e_G$. Its order divides $p$, so it is either $1$ or $p$.
It cannot be $1$, and hence $|\langle a\rangle|=p=|G|$. Therefore
$G=\langle a\rangle$.
\end{proof}

%\Needspace{12\baselineskip}
%\section{Towards classifying groups of order six}
%
%Let $|G|=6$. Lagrange's theorem shows that an element can have order
%$1$, $2$, $3$, or $6$. We first establish that elements of orders $2$ and
%$3$ must occur.
%
%\begin{lemma}
%Every finite group of even order contains an element of order $2$.
%\end{lemma}
%\begin{proof}
%Pair each element $a$ with its inverse $a^{-1}$. The distinct sets
%$\{a,a^{-1}\}$ partition $G$ into parts of size one or two. A part has
%size one exactly when $a=a^{-1}$.
%If the identity were the only such element, then $|G|$ would be one plus
%an even number, hence odd. Thus there is some $a\neq e_G$ with
%$a=a^{-1}$. It satisfies $a^2=e_G$ and has order $2$.
%\end{proof}
%
%\begin{lemma}
%If $a^2=e_G$ for every $a\in G$, then $G$ is abelian.
%\end{lemma}
%\begin{proof}
%Every element is its own inverse, so for $a,b\in G$,
%\[
%  ba=b^{-1}a^{-1}=(ab)^{-1}=ab.
%\]
%\end{proof}
%
%\begin{proposition}
%A group of order $6$ contains an element of order $3$.
%\end{proposition}
%\begin{proof}
%Suppose first that every nonidentity element has order $2$.
%The preceding lemma makes $G$ abelian. Choose distinct nonidentity
%elements $a,b$. Then
%\[
%  K=\{e_G,a,b,ab\}=\langle a,b\rangle
%\]
%is a subgroup: since $a,b$ commute and $a^2=b^2=e_G$, every product of
%powers of $a,b$ reduces to one of these four elements. They are distinct:
%$ab=e_G$ would imply $a=b$, while $ab=a$ or $ab=b$ would force $b=e_G$
%or $a=e_G$, respectively. Thus $|K|=4$, contradicting Lagrange's theorem
%because $4$ does not divide $6$.
%
%There is therefore an element $a$ whose order is neither $1$ nor $2$.
%Its order is $3$ or $6$. In the first case, take $x=a$.
%In the second case, take $x=a^2$; then $x^3=e_G$ and $x\neq e_G$, so
%$x$ has order $3$.
%\end{proof}
%
%Choose $x$ of order $3$ and $y$ of order $2$, and put
%$H=\langle x\rangle=\{e_G,x,x^2\}$. Since the elements of $H$ have
%orders $1$ or $3$, we have $y\notin H$. The coset criterion gives
%$yH\neq H$. Both cosets have three elements, so
%\begin{equation}
%  G=H\sqcup yH=\{e_G,x,x^2,y,yx,yx^2\}.
%\end{equation}
%In particular, $x$ and $y$ generate $G$.
%Determining how these generators multiply will be the next step in the
%classification.

\end{document}
